% ============================================================
%  Kapitel 5: Worker-Klassen und gemeinsame Infrastruktur
% ============================================================
\chapter{Worker-Klassen und gemeinsame Infrastruktur}

% ============================================================
\section{Funktionsweise}
% ============================================================

TilVista fuehrt alle datei-intensiven Operationen in
Hintergrund-Threads aus, damit der GUI-Thread jederzeit
reaktiv bleibt. Alle Worker folgen demselben Muster:

\begin{enumerate}[noitemsep]
  \item \klasse{QObject}-Subklasse mit einem \verb|run()|-Slot
  \item \verb|moveToThread()| in einen \klasse{QThread}
  \item Arbeit in \methode{run}, Ergebnis per Signal
  \item Selbstaufraeuumung via \verb|deleteLater()|
\end{enumerate}

Die Worker-Klassen im Ueberblick:
\begin{center}
\small
\begin{tabular}{lll}
\toprule
\textbf{Klasse} & \textbf{Aufgabe} & \textbf{Genutzt von} \\
\midrule
\klasse{ScanWorker}       & Verzeichnis rekursiv scannen  &
  AleaVueTab, MadoludusTab \\
\klasse{DBSaveWorker}     & JSON-Datei schreiben          &
  DirDatabasePanel, ShujukoPanel \\
\klasse{DBLoadWorker}     & JSON-Datei lesen              &
  DirDatabasePanel, ShujukoPanel \\
\klasse{CatalogueWriteWorker} & .dshow/.catalogue schreiben &
  DirDatabasePanel \\
\klasse{CatalogueLoadWorker}  & .dshow/.catalogue lesen     &
  DirDatabasePanel \\
\klasse{VideoFramesWorker}    & ffmpeg-Frames extrahieren   &
  VideoPreviewWidget, MadoludusWindow \\
\bottomrule
\end{tabular}
\end{center}

% ============================================================
\section{File by File}
% ============================================================

\subsection{scanworker.h / scanworker.cpp}

\begin{lstlisting}[style=header, caption={scanworker.h}]
class ScanWorker : public QObject {
    Q_OBJECT
public:
    explicit ScanWorker(const QString& directory,
                        QObject* parent = nullptr);
public slots:
    void run();
signals:
    void progressChanged(int percent);
    void resultReady(bool ok,
                     QStringList imageFiles,
                     QStringList allFiles);
private:
    QString m_directory;
};
\end{lstlisting}

\subsubsection{C++-Analyse: run() -- rekursiver Verzeichnis-Scan}

\begin{lstlisting}[caption={ScanWorker::run -- QDirIterator mit Subdirectories}]
void ScanWorker::run() {
    QStringList all, img;
    try {
        // Rekursive Iteration ohne Limit
        QDirIterator it(
            m_directory,
            QDir::Files | QDir::NoDotAndDotDot,
            QDirIterator::Subdirectories);

        while (it.hasNext()) {
            it.next();
            all << it.filePath();
        }
        const int total = all.isEmpty() ? 1 : all.size();
        const auto& suf = TV::imageSuffixes();

        for (int i = 0; i < all.size(); ++i) {
            const QString ext =
                '.' + QFileInfo(all[i]).suffix().toLower();
            if (suf.contains(ext))
                img << all[i];
            // Fortschritt in den GUI-Thread senden
            emit progressChanged((i+1)*100/total);
        }
        emit resultReady(true, img, all);
    } catch (...) {
        // Jede Exception abfangen: im Worker-Thread gibt
        // es keinen Aufrufer, der Exceptions behandeln koennte
        emit resultReady(false, {}, {});
    }
}
\end{lstlisting}

\begin{qtinfo}[title={Qt-Hintergrund: QDirIterator}]
\klasse{QDirIterator} traversiert einen Verzeichnisbaum
iterativ ohne Rekursion -- kein Stack-Overflow-Risiko auch
bei tief verschachtelten Strukturen.
Die Option \verb|Subdirectories| aktiviert die rekursive
Traversal. \verb|NoDotAndDotDot| schliesst \datei{.} und
\datei{..} aus.\\
Referenz: \url{https://doc.qt.io/qt-6/qdiriterator.html}
\end{qtinfo}

\paragraph{try/catch(...) im Worker-Thread}
Der \verb|catch(...)| (drei Punkte = jede Exception)
ist hier bewusst breit: Im Worker-Thread gibt es keinen
uebergeordneten Aufrufer auf dem Call-Stack, der eine
Exception sinnvoll behandeln koennte. Unbehandelte
Exceptions in einem Thread beenden das Programm unkontrolliert
(\verb|std::terminate|). Daher: abfangen, \verb|false|
melden, sauber beenden.

% ----------------------------------------------------------
\subsection{dbworkers.h / dbworkers.cpp}
% ----------------------------------------------------------

\begin{lstlisting}[style=header, caption={dbworkers.h}]
class DBSaveWorker : public QObject {
    Q_OBJECT
public:
    DBSaveWorker(const QString& path,
                 const QJsonObject& data,
                 QObject* p = nullptr);
public slots: void run();
signals: void resultReady(bool ok);
private:
    QString     m_path;
    QJsonObject m_data;  // Wert-Kopie (thread-sicher)
};

class DBLoadWorker : public QObject {
    Q_OBJECT
public:
    explicit DBLoadWorker(const QString& path,
                          QObject* p = nullptr);
public slots: void run();
signals: void resultReady(bool ok, QJsonObject data);
private:
    QString m_path;
};
\end{lstlisting}

\subsubsection{C++-Analyse: QJsonObject als Wert-Kopie}

\begin{lstlisting}[caption={DBSaveWorker::run}]
void DBSaveWorker::run() {
    QDir().mkpath(QFileInfo(m_path).absolutePath());
    QFile f(m_path);
    if (!f.open(QIODevice::WriteOnly | QIODevice::Truncate)) {
        emit resultReady(false);
        return;
    }
    // QJsonDocument::Indented: lesbare Formatierung
    f.write(QJsonDocument(m_data)
            .toJson(QJsonDocument::Indented));
    emit resultReady(true);
}
\end{lstlisting}

\member{m\_data} ist eine \emph{Wert-Kopie} von
\verb|QJsonObject|. Qt-Wertobjekte nutzen
\emph{Copy-on-Write (Implicit Sharing)}: Der Kopiervorgang
im Konstruktor ist sehr guenstig (nur Referenzzaehler
erhoehen). Erst wenn jemand schreibt, wird wirklich kopiert.
Ergebnis: Beide Threads (GUI und Worker) haben
unabhaengige Daten ohne Mutex.

\verb|QDir().mkpath(...)| erzeugt das Zielverzeichnis
rekursiv, falls es noch nicht existiert --
das Aequivalent zu \texttt{mkdir -p}.

% ----------------------------------------------------------
\subsection{catalogueworkers.h / catalogueworkers.cpp}
% ----------------------------------------------------------

\begin{lstlisting}[style=header, caption={catalogueworkers.h}]
class CatalogueWriteWorker : public QObject {
    Q_OBJECT
public:
    CatalogueWriteWorker(
        const QString& kaivoDir,
        const QString& dbPath,
        const QJsonObject& dbData,
        const QString& entryName,
        const QStringList& imageFiles,
        const QStringList& allFiles,
        QObject* p = nullptr);
public slots: void run();
signals: void resultReady(bool ok);
};

class CatalogueLoadWorker : public QObject {
    Q_OBJECT
public:
    CatalogueLoadWorker(
        const QString& kaivoDir,
        const QString& imgFname,
        const QString& allFname,
        QObject* p = nullptr);
public slots: void run();
signals:
    void resultReady(bool ok,
                     QStringList imageFiles,
                     QStringList allFiles);
private:
    QString m_kaivoDir, m_imgFname, m_allFname;
    // static: kein this-Pointer noetig
    static QStringList readLines(const QString& path);
};
\end{lstlisting}

\subsubsection{C++-Analyse: static readLines}

\verb|readLines| ist eine statische Hilfsmethode:
Sie braucht keinen Zustand des Objekts, nur einen Pfad.
\verb|static| im Klassenkontext bedeutet: kein impliziter
\verb|this|-Zeiger, kein Zugriff auf Instanz-Member.
Man haette auch eine freie Funktion nehmen koennen --
als statische Methode ist sie jedoch logisch der Klasse
zugeordnet.

\begin{lstlisting}[caption={CatalogueLoadWorker::readLines}]
QStringList CatalogueLoadWorker::readLines(
    const QString& path)
{
    QStringList r;
    QFile f(path);
    if (!f.open(QIODevice::ReadOnly | QIODevice::Text))
        return r;
    QTextStream in(&f);
    while (!in.atEnd()) {
        const QString l = in.readLine().trimmed();
        if (!l.isEmpty()) r << l;
    }
    return r;
}
\end{lstlisting}

\verb|QTextStream| liest zeilenweise und behandelt
Encoding-Fragen. \verb|trimmed()| entfernt fuehrende
und abschliessende Leerzeichen (inkl. \verb|\r| unter
Windows) -- wichtig fuer Cross-Platform-Kompatibilitaet
der Katalog-Dateien.

% ----------------------------------------------------------
\subsection{videoframesworker.h / videoframesworker.cpp}
% ----------------------------------------------------------

\klasse{VideoFramesWorker} extrahiert Einzelbilder aus
Videos via ffmpeg und legt sie als JPEG im Temp-Verzeichnis
ab.

\begin{lstlisting}[style=header, caption={videoframesworker.h -- Kern}]
class VideoFramesWorker : public QObject {
    Q_OBJECT
public:
    enum class Mode { IntervalSeconds, FixedCount };

    explicit VideoFramesWorker(
        const QString& videoPath,
        Mode    mode  = Mode::IntervalSeconds,
        double  value = 10.0,
        QObject* parent = nullptr);

    const QString& tmpDir() const { return m_tmpDir; }

public slots: void run();

signals:
    void resultReady(bool ok, QStringList framePaths);

private:
    // Videodauer via "ffmpeg -i" aus stderr lesen
    static double probeDurationSeconds(const QString& path);
    QList<double> computeTimestamps(double durationSeconds) const;

    QString m_videoPath, m_tmpDir;
    Mode    m_mode;
    double  m_value;
};
\end{lstlisting}

\subsubsection{C++-Analyse: Scoped Enum (enum class)}

\begin{lstlisting}
enum class Mode { IntervalSeconds, FixedCount };
\end{lstlisting}

\verb|enum class| (C++11) ist ein stark typisiertes Enum:
\begin{itemize}[noitemsep]
  \item Werte sind \emph{nicht} in den umgebenden Scope
    exportiert: Man schreibt
    \verb|Mode::IntervalSeconds|, nicht einfach
    \verb|IntervalSeconds|
  \item Kein implizites Konvertieren zu \verb|int|:
    Fehler werden vom Compiler erkannt statt still
    falsch zu laufen
  \item Besser lesbar als \verb|bool ffFixedCount|
    mit zwei Bedeutungen
\end{itemize}

\subsubsection{C++-Analyse: probeDurationSeconds}

\begin{lstlisting}[caption={probeDurationSeconds: ffmpeg-Stderr parsen}]
double VideoFramesWorker::probeDurationSeconds(
    const QString& path)
{
    // "ffmpeg -i <path>" gibt Metadaten auf stderr aus
    // und beendet sich mit Fehlercode (kein Output-File).
    // Genau das brauchen wir -- nur die Metadaten.
    QProcess proc;
    proc.start("ffmpeg", {"-i", path});
    proc.waitForFinished(8000);
    const QString out =
        QString::fromUtf8(proc.readAllStandardError());

    // "Duration: HH:MM:SS.xx" im stderr suchen
    static const QRegularExpression re(
        R"(Duration:\s*(\d+):(\d+):(\d+(?:\.\d+)?))");
    const auto m = re.match(out);
    if (!m.hasMatch()) return 0.0;

    return m.captured(1).toDouble() * 3600.0  // Stunden
         + m.captured(2).toDouble() *   60.0  // Minuten
         + m.captured(3).toDouble();           // Sekunden
}
\end{lstlisting}

\textbf{Warum stderr?} ffmpeg gibt Metadaten immer auf
\emph{stderr} aus, auch wenn es sich um normale
Informationen handelt. stdout ist fuer die eigentliche
Ausgabedatei reserviert.

\textbf{Raw-String-Regex}: \verb|R"(...)"| vermeidet
doppeltes Escaping. Ohne Raw-String:
\verb|"Duration:\\s*(\\d+):(\\d+):"| -- schlecht lesbar.

\subsubsection{v0.5.42: computeTimestamps -- durationssensitiv}

\begin{lstlisting}[caption={computeTimestamps: IntervalSeconds und FixedCount}]
QList<double> VideoFramesWorker::computeTimestamps(
    double durationSeconds) const
{
    constexpr int kMaxFrames = 60;
    QList<double> ts;

    if (durationSeconds <= 0.5) {
        ts << 0.0;
        return ts;
    }

    if (m_mode == Mode::IntervalSeconds) {
        // Ein Frame pro "value" Sekunden
        const double step = qMax(0.2, m_value);
        for (double t = 0.0; t < durationSeconds; t += step)
            ts << t;
        if (ts.isEmpty()) ts << 0.0;
    } else {
        // Genau "value" Frames gleichmaessig verteilt
        const int count = qMax(1, int(m_value));
        if (count == 1) {
            ts << durationSeconds / 2.0;
            return ts;
        }
        for (int i = 0; i < count; ++i)
            ts << durationSeconds * double(i)
                                  / double(count - 1);
    }

    // Auf 60 Frames deckeln (verhindert hunderte ffmpeg-Aufrufe)
    if (ts.size() > kMaxFrames) {
        QList<double> capped;
        for (int i = 0; i < kMaxFrames; ++i)
            capped << ts.at(
                i * (ts.size()-1) / (kMaxFrames-1));
        return capped;
    }
    return ts;
}
\end{lstlisting}

\begin{neuin}[title={Neu in v0.5.42: Dauer-bewusste Timestamps}]
Fruehre Version: fest \{0, 10, 20, 30, 40, 50\} Sekunden --
bei kurzen Videos (unter 10 Sekunden) gab es fast keine
Frames, bei langen Videos immer nur den Anfang.

Neue Version: Die Videodauer wird per \verb|ffmpeg -i|
ermittelt. Dann werden Timestamps relativ zur Gesamtdauer
berechnet -- fuer ein 2-Minuten-Video mit "1/sec" gibt
es 120 Frames (gedeckelt auf 60). Fuer "1/10 frames" gibt
es immer genau 10 gleichmaessig verteilte Frames.
\end{neuin}

\subsubsection{Qt-Analyse: QProcess fuer externe Programme}

\begin{lstlisting}[caption={run(): ffmpeg-Frames extrahieren}]
void VideoFramesWorker::run() {
    // ffmpeg vorhanden?
    if (QStandardPaths::findExecutable("ffmpeg").isEmpty()) {
        emit resultReady(false, {});
        return;
    }
    const double duration = probeDurationSeconds(m_videoPath);
    const QList<double> times = duration > 0.0
        ? computeTimestamps(duration)
        : QList<double>{0.0, 5.0, 10.0, 15.0, 20.0, 25.0};

    QStringList frames;
    int idx = 0;
    for (double t : times) {
        const QString out =
            QString("%1/f%2.jpg")
            .arg(m_tmpDir)
            // 3-stellig, Basis 10, Fuellzeichen '0'
            .arg(idx++, 3, 10, QChar('0'));

        QProcess proc;
        proc.start("ffmpeg", {
            "-i",        m_videoPath,
            "-ss",       QString::number(t, 'f', 2),
            "-vframes",  "1",
            "-f",        "image2",
            "-y",        out
        });
        proc.waitForFinished(12000);
        if (proc.exitCode() == 0 && QFile::exists(out))
            frames << out;
    }
    emit resultReady(!frames.isEmpty(), frames);
}
\end{lstlisting}

\paragraph{QString::arg mit Formatierung}
\verb|.arg(idx++, 3, 10, QChar('0'))| formatiert
eine Zahl mit:
\begin{itemize}[noitemsep]
  \item Breite 3 (\verb|3|)
  \item Basis 10 (\verb|10|)
  \item Fuellzeichen \verb|'0'|
\end{itemize}
Ergebnis: 0 $\to$ \verb|f000.jpg|, 10 $\to$ \verb|f010.jpg|.
Sortierbare Dateinamen ohne Extrabibliothek.

\begin{qtinfo}[title={Qt-Hintergrund: QProcess und Argumente}]
\klasse{QProcess} startet externe Programme sicher ohne
Shell-Interaktion. \verb|proc.start(programm, arglist)|
uebergibt Argumente als \klasse{QStringList} -- Leerzeichen
in Pfaden werden korrekt behandelt, Shell-Injection ist
nicht moeglich. Im Gegensatz zu
\verb|system("ffmpeg -i " + path)| ist das auch auf
Pfaden mit Leerzeichen und Sonderzeichen zuverlaessig.\\
Referenz: \url{https://doc.qt.io/qt-6/qprocess.html}
\end{qtinfo}

% ----------------------------------------------------------
\subsection{Tab 4 -- shortcutstab.h / shortcutstab.cpp}
% ----------------------------------------------------------

Der About-Tab ist ein reiner Informations-Tab ohne
Datenbankanbindung.

\subsubsection{Datentabellen als statische Structs}

\begin{lstlisting}[caption={Datenmodell als Compile-Zeit-Konstanten}]
struct Row     { const char* key; const char* description; };
struct Section { const char* title; QList<Row> rows; };

static const QList<Section> kSections = {
    {
        "Global",
        {
            {"Ctrl+Alt+F8", "Toggle Secret Mode"},
            {"Enter (DirBar)", "Verzeichnis uebernehmen"},
        }
    },
    {
        "AleaVue  --  Vollbild-Diashow",
        {
            {"->",      "Naechstes Bild (zufaellig)"},
            {"<-",      "Zurueck (Verlauf)"},
            {"S",       "Bookmark in shujuko"},
            {"Space",   "Timer pause/resume"},
            {"Esc",     "Diashow schliessen"},
        }
    },
    // ...
};
\end{lstlisting}

\verb|const char*| ist ein C-String-Zeiger direkt auf
das String-Literal im Programm-Daten-Segment. Da die
Strings nicht geaendert werden, ist das effizient und
sicher. \klasse{QList<Row>} kann mit
Initializer-Listen \verb|{ {...}, {...} }| befuellt werden.

\subsubsection{Dynamische Tabellenhoehe}

\begin{lstlisting}[caption={makeTable: exakte Hoehe berechnen}]
static QTableWidget* makeTable(const Section& sec) {
    auto* table = new QTableWidget(
        static_cast<int>(sec.rows.size()), 2);
    table->horizontalHeader()
         ->setStretchLastSection(true);
    table->verticalHeader()->setVisible(false);
    table->setEditTriggers(
        QAbstractItemView::NoEditTriggers);
    table->setSelectionMode(
        QAbstractItemView::NoSelection);
    table->setShowGrid(false);
    table->setAlternatingRowColors(true);

    // ... Eintraege setzen ...
    table->resizeRowsToContents();

    // Exakte Hoehe = Headerhoehe + alle Zeilenhoehen
    int h = table->horizontalHeader()->height();
    for (int r = 0; r < table->rowCount(); ++r)
        h += table->rowHeight(r);
    table->setFixedHeight(h + 4);  // +4 Pixel Puffer
    return table;
}
\end{lstlisting}

\verb|setFixedHeight| klemmt die Hoehe auf exakt den
Inhalt -- kein Leerraum, keine Scrollbar noetig.
\verb|setAlternatingRowColors(true)| erzeugt abwechselnde
Zeilenfaerbung automatisch -- per OS-Theme-Farben,
kein Hardcoding noetig.
